\part{Analyse}

%% ════════════════════════════════════════════════════════════════
\chapter{Fonctions num\'eriques}
\begin{prerequis}
Lire un repère, calculer la valeur d'une expression, résoudre graphiquement une
équation simple et distinguer abscisse et ordonnée.
\end{prerequis}
\begin{automatismes}
Pour $f(x)=x^2-3$, calculer $f(-2)$, $f(0)$ et $f(3)$; résoudre
$f(x)=1$; expliquer pourquoi $f(-2)=1$ ne signifie pas $f(1)=-2$.
\end{automatismes}
\begin{activite}[Activité d'entrée -- Une machine, plusieurs représentations]
Une machine associe à $x$ le nombre $x^2-4$. Construire une table pour sept
valeurs de $x$, placer les points correspondants, puis répondre aux mêmes
questions à partir de la formule, du tableau et du graphique. Comparer ce que
chaque représentation rend facile ou difficile.
\end{activite}


%% ════════════════════════════════════════════════════════════════

\begin{anecdote}[Euler et la notation $f(x)$]
Leonhard Euler ($1707$--$1783$) est le math\'ematicien le plus prolifique
de l'histoire : plus de $800$ publications. C'est lui qui introduisit la
notation $f(x)$ en $1734$. Il publia jusqu'\`a sa mort, m\^{e}me apr\`es \^{e}tre
devenu compl\`etement aveugle. Lorsqu'il perdit l'usage d'un \oe il, il aurait
dit : \og J'aurai moins de distractions.\fg
\end{anecdote}

\begin{objectifs}
\begin{itemize}
  \item Lire et interpr\'eter une courbe repr\'esentative.
  \item D\'eterminer domaine de d\'efinition, image, ant\'ec\'edent, extremums.
  \item Dresser et lire un tableau de variation.
  \item Conna\^itre les fonctions de r\'ef\'erence.
  \item D\'eduire les courbes des fonctions associ\'ees par transformations.
\end{itemize}
\end{objectifs}

\section{G\'en\'eralit\'es}

\begin{defn}[Fonction num\'erique]
Une \textbf{fonction num\'erique} $f$ est une application d'un ensemble $D_f \subset \R$
(domaine de d\'efinition) vers $\R$. $f(x)$ est l'\textbf{image} de $x$,
et $x$ est un \textbf{ant\'ec\'edent} de $f(x)$.
\end{defn}

\begin{defn}[Parit\'e]
$f$ est \textbf{paire} si $f(-x)=f(x)$ (courbe sym\'etrique par rapport \`a $Oy$).\\
$f$ est \textbf{impaire} si $f(-x)=-f(x)$ (courbe sym\'etrique par rapport \`a $O$).
\end{defn}

\section{Fonctions de r\'ef\'erence}

\begin{figure}[H]
\centering
\begin{tikzpicture}
  \begin{scope}
    \begin{axis}[width=6.3cm,height=5.2cm,axis equal image,axis lines=center,
      xlabel={$x$},ylabel={$y$},grid=both,major grid style={draw=BN!18},
      xmin=-2.2,xmax=2.2,ymin=-0.5,ymax=4.5,xtick distance=1,ytick distance=1,
      tick label style={font=\scriptsize},title style={font=\small\bfseries\color{BN}},
      title={$f(x)=x^2$}]
      \addplot[BN,very thick,domain=-2.1:2.1,samples=100]{x^2};
    \end{axis}
  \end{scope}
  \begin{scope}[xshift=7.4cm]
    \begin{axis}[width=6.3cm,height=5.2cm,axis equal image,axis lines=center,
      xlabel={$x$},ylabel={$y$},grid=both,major grid style={draw=VE!18},
      xmin=-0.2,xmax=4.2,ymin=-0.5,ymax=2.5,xtick distance=1,ytick distance=1,
      tick label style={font=\scriptsize},title style={font=\small\bfseries\color{VE}},
      title={$f(x)=\sqrt{x}$}]
      \addplot[VE,very thick,domain=0:4.1,samples=100]{sqrt(x)};
    \end{axis}
  \end{scope}
  \begin{scope}[yshift=-6.5cm]
    \begin{axis}[width=6.3cm,height=5.2cm,axis equal image,axis lines=center,
      xlabel={$x$},ylabel={$y$},grid=both,major grid style={draw=OR!18},
      xmin=-2.2,xmax=2.2,ymin=-0.5,ymax=2.5,xtick distance=1,ytick distance=1,
      tick label style={font=\scriptsize},title style={font=\small\bfseries\color{OR}},
      title={$f(x)=|x|$}]
      \addplot[OR,very thick,domain=-2.1:2.1,samples=100]{abs(x)};
    \end{axis}
  \end{scope}
  \begin{scope}[xshift=7.4cm,yshift=-6.5cm]
    \begin{axis}[width=6.3cm,height=5.2cm,axis equal image,axis lines=center,
      xlabel={$x$},ylabel={$y$},grid=both,major grid style={draw=VI!18},
      xmin=-1.7,xmax=1.7,ymin=-4.2,ymax=4.2,xtick distance=1,ytick distance=1,
      tick label style={font=\scriptsize},title style={font=\small\bfseries\color{VI}},
      title={$f(x)=x^3$}]
      \addplot[VI,very thick,domain=-1.6:1.6,samples=100]{x^3};
    \end{axis}
  \end{scope}
\end{tikzpicture}
\caption{Fonctions de référence au programme : graduations régulières et identiques sur chaque axe.}
\end{figure}

\horsprogramme{Approfondissement graphique sur les translations et les symétries de courbes.}
\section{Transformations}

\begin{prop}[Transformations \'el\'ementaires]
\`A partir de la courbe $\mathcal{C}_f$ de $f$ :
\begin{itemize}
  \item $f(x)+a$ : translation verticale de vecteur $(0,a)$.
  \item $f(x+a)$ : translation horizontale de vecteur $(-a,0)$.
  \item $-f(x)$ : sym\'etrie par rapport \`a $Ox$.
  \item $f(-x)$ : sym\'etrie par rapport \`a $Oy$.
  \item $|f(x)|$ : replier la partie n\'egative vers le haut.
\end{itemize}
\end{prop}


\section{Lecture de courbe et vocabulaire}

\begin{defn}[Domaine de d\'efinition]
\index{domaine de définition}
Le \textbf{domaine de d\'efinition} $D_f$ est l'ensemble des r\'eels $x$
pour lesquels $f(x)$ est d\'efini. Deux causes d'exclusion~:
\begin{itemize}
  \item Division par z\'ero~: exclure les $x$ qui annulent le d\'enominateur.
  \item Racine carr\'ee d'un n\'egatif~: exclure les $x$ o\`u l'expression
        sous le radical est $<0$.
\end{itemize}
\end{defn}

\begin{defn}[Image et ant\'ec\'edent]
\index{image}\index{antécédent}
Pour $y = f(x)$~:
\begin{itemize}
  \item $y$ est l'\textbf{image} de $x$ par $f$.
  \item $x$ est un \textbf{ant\'ec\'edent} de $y$~: il peut y en avoir
        plusieurs, un seul, ou aucun.
\end{itemize}
\emph{Graphiquement~:} l'image de $x=a$ est l'ordonn\'ee de la courbe
\`a l'abscisse $a$~; les ant\'ec\'edents de $y=b$ sont les abscisses
des points d'intersection de la courbe avec la droite horizontale $y=b$.
\end{defn}

\begin{defn}[Monotonie et extremums]
\index{monotonie}\index{extremum}
\begin{itemize}
  \item $f$ est \textbf{croissante} sur $I$ si~:
        $\forall x_1,x_2\in I,\ x_1<x_2\Rightarrow f(x_1)<f(x_2)$.
  \item $f$ est \textbf{d\'ecroissante} sur $I$ si~:
        $\forall x_1,x_2\in I,\ x_1<x_2\Rightarrow f(x_1)>f(x_2)$.
  \item Un \textbf{minimum} (resp.\ maximum) est une valeur de $f$ plus
        petite (resp.\ grande) que toutes les valeurs voisines.
\end{itemize}
\end{defn}

\begin{methode}[Dresser un tableau de variations]
\index{tableau de variations}
\begin{enumerate}
  \item Identifier le domaine de d\'efinition.
  \item D\'eterminer les intervalles de croissance et d\'ecroissance
        (graphiquement, ou en utilisant le signe de $f(x_2)-f(x_1)$).
  \item Indiquer les valeurs aux bornes et aux extremums.
\end{enumerate}
\end{methode}

\begin{ex}[Tableau de variations de $f(x)=x^2-4$]
\begin{center}
\renewcommand{\arraystretch}{1.8}
\begin{tabular}{|c|ccccc|}
\hline
$x$ & $-\infty$ & & $0$ & & $+\infty$ \\
\hline
$f(x)$ & $+\infty$ & \begin{tikzpicture}[baseline=-3pt]
  \draw[-latex,BO,thick] (0,0.3)--(0.8,-0.3);
\end{tikzpicture} & $-4$ & \begin{tikzpicture}[baseline=-3pt]
  \draw[-latex,VE,thick] (0,-0.3)--(0.8,0.3);
\end{tikzpicture} & $+\infty$ \\
\hline
\end{tabular}
\end{center}
$f$ est d\'ecroissante sur $]-\infty,0]$ et croissante sur $[0,+\infty[$~;
minimum $f(0)=-4$.
\end{ex}

\section{Fonctions de r\'ef\'erence}

\begin{center}
\renewcommand{\arraystretch}{1.5}
\begin{tabular}{|l|c|c|l|}
\hline
\rowcolor{BN!10}
Fonction & $D_f$ & Parit\'e & Variations \\
\hline
$f(x)=x^2$   & $\R$ & paire   & d\'ecr.\ sur $]-\infty,0]$, croiss.\ sur $[0,+\infty[$ \\
$f(x)=x^3$   & $\R$ & impaire & croissante sur $\R$ \\
$f(x)=\sqrt{x}$ & $[0,+\infty[$ & --- & croissante sur $[0,+\infty[$ \\
$f(x)=|x|$   & $\R$ & paire   & d\'ecr.\ sur $]-\infty,0]$, croiss.\ sur $[0,+\infty[$ \\
$f(x)=1/x$   & $\R^*$ & impaire & d\'ecroissante sur $]-\infty,0[$ et sur $]0,+\infty[$ \\
\hline
\end{tabular}
\end{center}

\begin{figure}[H]
\centering
\begin{tikzpicture}
  % 3 columns × 2 rows of small axes, side-by-side
  % Each function gets its own clean axes
  \begin{scope}[xshift=0cm]
    \begin{axis}[width=4.2cm,height=4.2cm,axis lines=center,
      xlabel={$x$},ylabel={$y$},
      xmin=-2.2,xmax=2.2,ymin=-0.4,ymax=4.8,
      xtick={-2,-1,0,1,2},ytick={0,2,4},
      tick label style={font=\tiny},
      title style={font=\small\bfseries\color{BN}},
      title={$f(x)=x^2$}]
      \addplot[BN,very thick,domain=-2.1:2.1,samples=80]{x^2};
      \node[BN,font=\tiny,anchor=west] at (axis cs:1.2,3.2) {$x^2$};
    \end{axis}
  \end{scope}
  \begin{scope}[xshift=4.5cm]
    \begin{axis}[width=4.2cm,height=4.2cm,axis lines=center,
      xlabel={$x$},ylabel={$y$},
      xmin=-0.3,xmax=4.5,ymin=-0.3,ymax=2.4,
      xtick={0,1,2,3,4},ytick={0,1,2},
      tick label style={font=\tiny},
      title style={font=\small\bfseries\color{VE}},
      title={$f(x)=\sqrt{x}$}]
      \addplot[VE,very thick,domain=0:4.4,samples=80]{sqrt(x)};
      \node[VE,font=\tiny,anchor=west] at (axis cs:2.5,1.65) {$\sqrt{x}$};
    \end{axis}
  \end{scope}
  \begin{scope}[xshift=9cm]
    \begin{axis}[width=4.2cm,height=4.2cm,axis lines=center,
      xlabel={$x$},ylabel={$y$},
      xmin=-2.5,xmax=2.5,ymin=-0.3,ymax=2.8,
      xtick={-2,-1,0,1,2},ytick={0,1,2},
      tick label style={font=\tiny},
      title style={font=\small\bfseries\color{OR}},
      title={$f(x)=|x|$}]
      \addplot[OR,very thick,domain=-2.4:2.4,samples=80]{abs(x)};
      \node[OR,font=\tiny,anchor=west] at (axis cs:0.5,1.2) {$|x|$};
    \end{axis}
  \end{scope}
  \begin{scope}[xshift=0cm,yshift=-4.8cm]
    \begin{axis}[width=4.2cm,height=4.2cm,axis lines=center,
      xlabel={$x$},ylabel={$y$},
      xmin=-2.5,xmax=2.5,ymin=-4.5,ymax=4.5,
      xtick={-2,-1,0,1,2},ytick={-4,-2,0,2,4},
      tick label style={font=\tiny},
      title style={font=\small\bfseries\color{BO}},
      title={$f(x)=1/x$}]
      \addplot[BO,very thick,domain=0.22:2.4,samples=80]{1/x};
      \addplot[BO,very thick,domain=-2.4:-0.22,samples=80]{1/x};
      \node[BO,font=\tiny,anchor=west] at (axis cs:0.8,2.5) {$1/x$};
    \end{axis}
  \end{scope}
  \begin{scope}[xshift=4.5cm,yshift=-4.8cm]
    \begin{axis}[width=4.2cm,height=4.2cm,axis lines=center,
      xlabel={$x$},ylabel={$y$},
      xmin=-1.8,xmax=1.8,ymin=-4.5,ymax=4.5,
      xtick={-1,0,1},ytick={-4,-2,0,2,4},
      tick label style={font=\tiny},
      title style={font=\small\bfseries\color{VI}},
      title={$f(x)=x^3$}]
      \addplot[VI,very thick,domain=-1.65:1.65,samples=80]{x^3};
      \node[VI,font=\tiny,anchor=west] at (axis cs:0.6,1.8) {$x^3$};
    \end{axis}
  \end{scope}
  % 6th cell: summary table of key properties
  \begin{scope}[xshift=9cm,yshift=-4.8cm]
    \node[draw=BN!30,fill=GC,rounded corners=3pt,
          inner sep=6pt,text width=3.6cm,font=\tiny,
          align=left] at (1.8,1.4) {
      \textbf{R\'ecap\'itulatif}\\[2pt]
      $x^2$~: min en $x=0$, paire\\
      $\sqrt{x}$~: croissante sur $[0,+\infty[$\\
      $|x|$~: min en $x=0$, paire\\
      $1/x$~: d\'ecr. sur $\R^*$, impaire\\
      $x^3$~: croissante sur $\R$, impaire
    };
  \end{scope}
\end{tikzpicture}
\caption{Les cinq fonctions de r\'ef\'erence --- chacune avec son domaine et ses propri\'et\'es.}
\end{figure}

\section{Transformations de courbes}


\begin{prop}[Transformations \'el\'ementaires]
\index{transformations de courbes}
\`A partir de la courbe $\mathcal{C}_f$ de $f$~:
\begin{center}
\small
\renewcommand{\arraystretch}{1.4}
\begin{tabularx}{\linewidth}{|l|l|X|}
\hline
\rowcolor{BN!10} Nouvelle fonction & Transformation & Description \\
\hline
$f(x)+k$ & Translation verticale & Monte de $k$ (ou descend si $k<0$) \\
$f(x+k)$ & Translation horizontale & Glisse \`a gauche de $k$ (si $k>0$) \\
$-f(x)$ & Sym\'etrie par rapport \`a $Ox$ & Retourne la courbe \\
$f(-x)$ & Sym\'etrie par rapport \`a $Oy$ & Image miroir \\
$|f(x)|$ & Valeur absolue & Remonte la partie n\'egative \\
\hline
\end{tabularx}
\end{center}
\end{prop}

\begin{ex}[Transformations de $x^2$]
\begin{itemize}
  \item $x^2+3$~: parabole mont\'ee de $3$ (sommet en $(0,3)$).
  \item $(x-2)^2$~: parabole d\'eplac\'ee \`a droite de $2$ (sommet en $(2,0)$).
  \item $-(x+1)^2+2$~: parabole retourn\'ee, sommet en $(-1,2)$, maximum $2$.
\end{itemize}
\end{ex}

\section{Exercices}

\subsection*{\diff{1}\quad Niveau 1 \textemdash\ Vocabulaire et fonctions de r\'ef\'erence}

\begin{exo}[Domaine de d\'efinition \corrige]{1}
D\'eterminer le domaine de d\'efinition~:
\begin{enumerate}
  \item $f(x) = \sqrt{2x-4}$
  \item $g(x) = \dfrac{1}{x^2-1}$
  \item $h(x) = \dfrac{\sqrt{x+3}}{x-2}$
  \item $k(x) = \sqrt{9-x^2}$
  \item $p(x) = \dfrac{1}{\sqrt{x-1}}$
  \item $q(x) = \sqrt{x^2-4}$
  \item $r(x) = \dfrac{x+1}{\sqrt{3-x}}$
  \item $s(x) = \sqrt{x} + \sqrt{2-x}$
\end{enumerate}
\end{exo}

\begin{exo}[Lire une courbe \corrige]{1}
\`A partir de la courbe repr\'esentative de $f(x)=x^2-2x-3$~:
\begin{enumerate}
  \item Dresser le tableau de valeurs pour $x \in \{-2,-1,0,1,2,3,4\}$.
  \item Tracer la courbe.
  \item Lire graphiquement~: le minimum de $f$, les z\'eros de $f$,
        l'image de $x=3$, les ant\'ec\'edents de $y=5$.
  \item V\'erifier les z\'eros et le minimum par le calcul.
\end{enumerate}
\end{exo}

\begin{exo}[Parit\'e \corrige]{1}
\'Etudier la parit\'e (paire, impaire, ou ni l'une ni l'autre)~:
\begin{multicols}{2}\small
\begin{enumerate}
  \item $f(x) = x^4 - 3x^2$
  \item $g(x) = x^3 + x$
  \item $h(x) = x^2 + x$
  \item $p(x) = x^4 - x^2 + 1$
  \item $q(x) = x^3 - x$
  \item $r(x) = |x| + 1$
  \item $s(x) = x^2 - 2x$
  \item $t(x) = \dfrac{x^2}{x^2+1}$
\end{enumerate}
\end{multicols}
\end{exo}

\begin{exo}[Tableau de variations --- lecture \corrige]{1}
On donne le tableau de variations d'une fonction $f$~:
\begin{center}
\renewcommand{\arraystretch}{1.8}
\begin{tabular}{|c|ccccccc|}
\hline
$x$ & $-3$ & & $-1$ & & $2$ & & $5$ \\
\hline
$f(x)$ & $1$ & $\nearrow$ & $4$ & $\searrow$ & $-2$ & $\nearrow$ & $3$ \\
\hline
\end{tabular}
\end{center}
\begin{enumerate}
  \item Sur quels intervalles $f$ est-elle croissante~? D\'ecroissante~?
  \item Quel est le maximum de $f$~? Le minimum~?
  \item Peut-on affirmer que $f(-2)=0$~? Que $f(0)=1$~?
  \item Combien d'ant\'ec\'edents de $0$ peut-il y avoir~?
\end{enumerate}
\end{exo}

\begin{exo}[Fonctions de r\'ef\'erence~: valeurs \corrige]{1}
Sans calculatrice, calculer les valeurs suivantes et placer les points
sur un graphe~:
\begin{multicols}{2}
\begin{enumerate}
  \item $f(x)=x^2$ pour\\ $x\in\{-3,-2,-1,0,1,2,3\}$
  \item $g(x)=\sqrt{x}$ pour $x\in\{0,1,4,9,16\}$
  \item $h(x)=1/x$ pour $x\in\{-2,-1,\frac{1}{2},1,2\}$
  \item $k(x)=x^3$ pour $x\in\{-2,-1,0,1,2\}$
\end{enumerate}
\end{multicols}
\end{exo}

\begin{exo}[Asymptotes \corrige]{1}
\index{asymptote}
Pour chaque fonction, donner le domaine et indiquer les asymptotes~:
\begin{enumerate}
  \item $f(x) = \dfrac{1}{x}$~: tracer et indiquer les asymptotes $x=0$ et $y=0$.
  \item $g(x) = \dfrac{1}{x} + 2$~: quelle est la nouvelle asymptote horizontale~?
  \item $h(x) = \dfrac{1}{x-1}$~: identifier les asymptotes.
  \item $k(x) = \dfrac{1}{x+3} - 1$~: identifier les asymptotes.
\end{enumerate}
\end{exo}

\begin{exo}[Image et ant\'ec\'edent \corrige]{1}
Pour $f(x) = x^2 - 4$~:
\begin{enumerate}
  \item Calculer $f(0)$, $f(2)$, $f(-3)$, $f(\sqrt{5})$.
  \item Trouver les ant\'ec\'edents de $0$, de $5$, de $-5$.
  \item $f$ a-t-elle toujours deux ant\'ec\'edents~? Un seul~? Aucun~?
  \item Tracer la courbe et v\'erifier graphiquement.
\end{enumerate}
\end{exo}

\begin{exo}[Comparaison de fonctions \corrige]{1}
Comparer $f$ et $g$ sur les intervalles indiqu\'es~:
\begin{enumerate}
  \item $f(x)=x^2$ et $g(x)=x$ sur $[0,1]$ puis sur $[1,2]$.
        Tracer et d\'eterminer graphiquement o\`u $f(x)\geq g(x)$.
  \item $f(x)=\sqrt{x}$ et $g(x)=x^2$ sur $[0,1]$ puis sur $[1,4]$.
  \item $f(x)=x$ et $g(x)=1/x$ sur $]0,1[$ puis sur $]1,+\infty[$.
\end{enumerate}
\end{exo}

\begin{exo}[Tableau de variations de fonctions de r\'ef\'erence \corrige]{1}
Dresser le tableau de variations complet~:
\begin{multicols}{2}
\begin{enumerate}
  \item $f(x) = x^2 - 4$
  \item $g(x) = -x^2 + 2x$
  \item $h(x) = x^2 - 6x + 5$
  \item $k(x) = -x^3$
  \item $p(x) = \sqrt{x+4}$
  \item $q(x) = \dfrac{1}{x+2}$
\end{enumerate}
\end{multicols}
\end{exo}

\begin{exo}[Transformations~: lire la nouvelle courbe \corrige]{1}
\`A partir de la courbe de $f(x)=x^2$, d\'ecrire et tracer~:
\begin{multicols}{2}
\begin{enumerate}
  \item $g(x) = x^2 + 3$
  \item $h(x) = x^2 - 2$
  \item $k(x) = (x-2)^2$
  \item $p(x) = (x+1)^2$
  \item $q(x) = -x^2$
  \item $r(x) = -(x-1)^2 + 3$
\end{enumerate}
\end{multicols}
\end{exo}

\begin{exo}[Variations et sens de variation \corrige]{1}
En utilisant les valeurs du tableau, d\'eterminer le sens de variation~:
\begin{enumerate}
  \item $f(x)=3x+1$~: calculer $f(0)$, $f(1)$, $f(2)$. $f$ est-elle croissante~?
  \item $g(x)=-2x+5$~: calculer $g(0)$, $g(1)$, $g(2)$.
        $g$ est-elle d\'ecroissante~?
  \item $h(x)=x^2-2x$~: calculer $h(-1)$, $h(0)$, $h(1)$, $h(2)$, $h(3)$.
        Sur quel(s) intervalle(s) $h$ semble-t-elle croissante~? D\'ecroissante~?
\end{enumerate}
\end{exo}

\begin{exo}[Extremum et minimum \corrige]{1}
\begin{enumerate}
  \item La fonction $f(x)=x^2-4x+7$ est-elle toujours positive~?
        Mettre sous forme canonique pour trouver le minimum.
  \item La fonction $g(x)=-x^2+6x-8$~: trouver le maximum par
        mise sous forme canonique.
  \item Pour $h(x)=2x^2-8x+3$~: trouver le minimum et les z\'eros.
  \item Pour $k(x)=|x-3|$~: trouver le minimum. En quel point est-il atteint~?
\end{enumerate}
\end{exo}

\begin{exo}[Fonctions affines \corrige]{1}
\index{fonction affine}
Une \textbf{fonction affine} est de la forme $f(x)=mx+p$ ($m\ne0$)~;
si $m=0$, c'est une \textbf{fonction constante}.
\begin{enumerate}
  \item Calculer la pente et l'ordonn\'ee \`a l'origine de~:
        $f(x)=3x-1$~; $g(x)=-2x+4$~; $h(x)=\frac{x}{2}+3$.
  \item Quelle fonction affine passe par $A(0,2)$ et $B(3,8)$~?
  \item Quelle fonction affine est croissante, a un z\'ero en $x=4$
        et $f(0)=6$~?
  \item La droite $y=mx+p$ croise l'axe $Ox$ en $-p/m$ (si $m\ne0$).
        V\'erifier pour les fonctions de la question~(1).
\end{enumerate}
\end{exo}

\begin{exo}[Compos\'ee de transformations \corrige]{1}
D\'ecrire la transformation qui passe de $f$ \`a $g$~:
\begin{enumerate}
  \item $f(x)=\sqrt{x}$ et $g(x)=\sqrt{x+4}$
  \item $f(x)=1/x$ et $g(x)=1/x-3$
  \item $f(x)=x^3$ et $g(x)=-(x-1)^3$
  \item $f(x)=x^2$ et $g(x)=|x^2-1|$
\end{enumerate}
\end{exo}

\begin{exo}[Valeur absolue : tracer et r\'esoudre \corrige]{1}
\begin{enumerate}
  \item Tracer $f(x)=|x-2|$~: identifier les intervalles de croissance.
  \item Tracer $g(x)=|x+1|-2$.
  \item R\'esoudre graphiquement puis alg\'ebriquement $|x-2|=3$.
  \item R\'esoudre $|x+1|-2=1$.
\end{enumerate}
\end{exo}

\subsection*{\diff{2}\quad Niveau 2 \textemdash\ Applications et \'etude de fonctions}

\begin{exo}[\'Etude compl\`ete d'une fonction \corrige]{2}
Pour $f(x)=\sqrt{x}$ sur $[0,9]$~:
\begin{enumerate}
  \item D\'eterminer le domaine et les ant\'ec\'edents de $2$, $3$ et $4$.
  \item Dresser le tableau de variations.
  \item \'Etudier la position relative des courbes $y=\sqrt{x}$ et $y=x/3$
        en r\'esolvant $\sqrt{x}=x/3$.
  \item Tracer les deux courbes sur $[0,9]$.
\end{enumerate}
\end{exo}

\begin{exo}[Valeur absolue et \'etude \corrige]{2}
Pour $f(x)=|2x-4|$~:
\begin{enumerate}
  \item \'Ecrire $f$ sans valeur absolue (par cas~: $x<2$ et $x\geq2$).
  \item Dresser le tableau de variations.
  \item R\'esoudre graphiquement puis alg\'ebriquement $f(x)=x$.
  \item \'Etudier le signe de $f(x)-x$.
\end{enumerate}
\end{exo}

\begin{exo}[Variations d'une fonction du second degr\'e \corrige]{2}
\'Etudier les variations de $f(x)=ax^2+bx+c$ en utilisant la forme canonique~:
\begin{enumerate}
  \item $f(x)=x^2-4x+1$~: trouver le sommet, les intervalles de variation.
  \item $g(x)=-2x^2+8x-3$~: trouver le sommet et le maximum.
  \item D\'emontrer que si $a>0$, $f$ est d\'ecroissante sur $]-\infty,\alpha]$
        et croissante sur $[\alpha,+\infty[$ o\`u $\alpha=-\dfrac{b}{2a}$.
        \emph{(Comparer $f(x_2)-f(x_1)$ pour $x_1<x_2$, avec ou sans $\alpha$.)}
\end{enumerate}
\end{exo}

\begin{exo}[Fonctions et in\'egalit\'es \corrige]{2}
\begin{enumerate}
  \item Montrer que $x^2\geq x$ pour $x\geq1$ et $x^2\leq x$ pour $0\leq x\leq1$.
        \emph{(\'Etudier le signe de $x^2-x=x(x-1)$.)}
  \item Montrer que $\sqrt{x}\leq x$ pour $x\geq1$ et $\sqrt{x}\geq x$ pour $0\leq x\leq1$.
        \emph{(\'Etudier $x^2\leq x$ en \'echangeant $x$ et $\sqrt{x}$.)}
  \item Pour quelles valeurs de $x>0$ a-t-on $1/x\geq x$~?
  \item Tracer les quatre fonctions sur un m\^eme graphe et v\'erifier visuellement.
\end{enumerate}
\end{exo}

\begin{exo}[Compos\'ee de fonctions de r\'ef\'erence \corrige]{2}
Pour chaque paire de fonctions $f$ et $g$, calculer $f(g(x))$ et donner son domaine~:
\begin{multicols}{2}
\begin{enumerate}
  \item $f(x)=x^2$, $g(x)=x-1$
  \item $f(x)=\sqrt{x}$, $g(x)=4-x$
  \item $f(x)=1/x$, $g(x)=x^2$
  \item $f(x)=|x|$, $g(x)=x^2-1$
  \item $f(x)=\sqrt{x}$, $g(x)=x^2$
  \item $f(x)=x^2$, $g(x)=\sqrt{x}$
\end{enumerate}
\end{multicols}
\end{exo}

\begin{exo}[Position relative de deux courbes \corrige]{2}
\'Etudier la position relative des courbes de $f$ et $g$, puis tracer~:
\begin{enumerate}
  \item $f(x)=x^2-1$ et $g(x)=x+1$.
  \item $f(x)=\sqrt{x}$ et $g(x)=x/2+1$ sur $[0,9]$.
  \item $f(x)=1/x$ et $g(x)=4-x$ sur $]0,+\infty[$.
\end{enumerate}
\end{exo}

\begin{exo}[Monotonie et in\'egalit\'es \corrige]{2}
\begin{enumerate}
  \item Soit $f$ croissante sur $\R$. Si $f(3)=5$ et $f(7)=9$, que peut-on
        dire de $f(5)$~? De $f(x)$ pour $x<3$~?
  \item Soit $g$ d\'ecroissante sur $[0,4]$, avec $g(0)=10$ et $g(4)=2$.
        Encadrer $g(1)$, $g(2)$, $g(3)$.
  \item D\'emontrer que $f(x)=x^3$ est croissante sur $\R$~:
        pour $x_1<x_2$, calculer $f(x_2)-f(x_1)=x_2^3-x_1^3$
        et montrer que c'est positif.
        \emph{(Factoriser $a^3-b^3=(a-b)(a^2+ab+b^2)$ et \'etudier
        le signe de chaque facteur.)}
\end{enumerate}
\end{exo}

\begin{exo}[Application~: trajectoire \corrige]{2}
Un objet est lanc\'e verticalement. Sa hauteur (en m) \`a l'instant $t$ (en s)
est $h(t) = -5t^2 + 20t + 1$.
\begin{enumerate}
  \item Calculer $h(0)$, $h(1)$, $h(2)$, $h(3)$, $h(4)$.
  \item Dresser le tableau de variations de $h$.
  \item Quelle est la hauteur maximale atteinte, et quand~?
  \item \`A quel instant l'objet est-il \`a hauteur $16$~m~?
  \item Quand retouche-t-il le sol ($h=0$)~?
\end{enumerate}
\end{exo}

\begin{exo}[Fonctions et racines \corrige]{2}
Pour $f(x)=x^2-2$~:
\begin{enumerate}
  \item Montrer que $f$ s'annule pour $x=\sqrt{2}$ et $x=-\sqrt{2}$.
  \item Dresser le tableau de signe de $f(x)$.
  \item La fonction $g(x)=\sqrt{x^2-2}$ est-elle d\'efinie sur $\R$~?
        Donner son domaine.
  \item Pour $h(x)=\dfrac{1}{x^2-2}$~: donner le domaine et les
        \'equations des asymptotes verticales.
\end{enumerate}
\end{exo}

\begin{exo}[Fonctions et algorithme \computer \corrige]{2}
\'Ecrire un programme Python qui~:
\begin{enumerate}
  \item Calcule et affiche un tableau de valeurs de $f(x)=x^2-3x+1$ pour
        \[x\in\{-2,-1,0,1,2,3,4\}.\]
  \item Trouve num\'eriquement le minimum de $f$ sur $[-2,4]$ en testant
        les valeurs avec un pas de $0{,}1$.
  \item V\'erifier analytiquement en utilisant la forme canonique.
\end{enumerate}
\end{exo}

\begin{exo}[Tableau de variations~: compl\'eter \corrige]{2}
Pour chaque fonction, compl\'eter le tableau de variations et indiquer
le (ou les) extremums~:
\begin{enumerate}
  \item $f(x)=x^2-6x+5$ sur $[-1,5]$
  \item $g(x)=\sqrt{x-1}$ sur $[1,10]$
  \item $h(x)=|3x-6|$ sur $[-2,6]$
  \item $k(x)=\dfrac{1}{x-2}$ sur $]2,+\infty[$
\end{enumerate}
\end{exo}

\begin{exo}[Encadrement par monotonie \corrige]{2}
Soit $f$ croissante sur $[a,b]$. On sait que $f(2)=5$ et $f(6)=9$.
\begin{enumerate}
  \item Peut-on affirmer que $5 \leq f(4) \leq 9$~?
  \item Peut-on affirmer que $f(x) \geq 5$ pour tout $x\in[2,6]$~?
  \item Pour $f(x)=\sqrt{x}$, encadrer $\sqrt{5}$ entre $f(4)$ et $f(9)$.
  \item Affiner~: encadrer $\sqrt{5}$ entre $f(4)$ et $f(4{,}84)$
        sachant que $4{,}84=2{,}2^2$.
\end{enumerate}
\end{exo}

\begin{exo}[Fonctions et r\'esolution graphique \corrige]{2}
Pour $f(x)=x^2-2x$ et $g(x)=2x-3$~:
\begin{enumerate}
  \item Tracer les deux courbes sur $[-1,4]$.
  \item Lire graphiquement le(s) point(s) d'intersection.
  \item R\'esoudre alg\'ebriquement $f(x)=g(x)$.
  \item R\'esoudre $f(x)\leq g(x)$ et repr\'esenter sur un axe.
\end{enumerate}
\end{exo}

\begin{exo}[Domaine de fonctions compos\'ees \corrige]{2}
Donner le domaine de d\'efinition et simplifier si possible~:
\begin{multicols}{2}
\begin{enumerate}
  \item $f(x)=\sqrt{x^2-3x+2}$
  \item $g(x)=\dfrac{1}{\sqrt{1-x^2}}$
  \item $h(x)=\sqrt{\dfrac{x-1}{x+1}}$
  \item $k(x)=\dfrac{\sqrt{x}}{\sqrt{x}-1}$
\end{enumerate}
\end{multicols}
\end{exo}

\begin{exo}[Lire un tableau de variations \corrige]{1}
Pour une fonction $f$ dont le tableau de variations est~:
\begin{center}
\renewcommand{\arraystretch}{1.8}
\begin{tabular}{|c|ccccccc|}
\hline
$x$ & $-2$ & & $0$ & & $3$ & & $5$ \\
\hline
$f$ & $-1$ & $\nearrow$ & $4$ & $\searrow$ & $0$ & $\nearrow$ & $2$ \\
\hline
\end{tabular}
\end{center}
\begin{enumerate}
  \item Combien de solutions a l'\'equation $f(x)=2$~?
  \item Combien de solutions a $f(x)=1$~? $f(x)=5$~? $f(x)=-2$~?
  \item Sur quel intervalle $f(x)\geq0$~?
  \item Quel est l'extremum absolu de $f$ sur $[-2,5]$~?
\end{enumerate}
\end{exo}

\begin{exo}[Fonctions et tableur \computer \corrige]{2}
\begin{enumerate}
  \item \'Ecrire un programme Python qui calcule le tableau de valeurs
        de $f(x)=x^3-3x$ pour $x\in\{-3,-2,-1,0,1,2,3\}$.
  \item Identifier num\'eriquement les valeurs de $x$ o\`u $f$ semble changer
        de monotonie.
  \item D\'eterminer graphiquement les z\'eros de $f$.
        Les v\'erifier alg\'ebriquement ($f(x)=x(x^2-3)$).
\end{enumerate}
\end{exo}

\begin{exo}[\logicoalgo\ Algorithme~: chercher les z\'eros de $f$ \corrige]{2}

\begin{algorithm}[H]
\caption{Recherche de z\'ero par balayage}
\KwIn{une fonction $f$, une borne inf. $a$, une borne sup. $b$, un pas $h$}
\KwOut{les $x$ o\`u $f$ semble changer de signe}
$x \leftarrow a$\;
\TantQue{$x + h \leq b$}{
  \Si{$f(x) \times f(x+h) < 0$}{
    \Afficher \og Changement de signe dans $[$\fg, $x$, \og$;$\fg, $x+h$, \og$]$\fg\;
  }
  $x \leftarrow x + h$\;
}
\end{algorithm}

\begin{enumerate}
  \item La condition $f(x)\times f(x+h)<0$ est-elle une condition n\'ecessaire
        ou suffisante pour qu'il y ait une racine dans $[x,x+h]$~?
        Donner un contre-exemple si elle n'est que suffisante.
  \item Appliquer (sur papier, avec $h=1$) \`a $f(x)=x^2-5$ sur $[-3,3]$~:
        quel(s) intervalle(s) contiennent une racine~?
  \item Combien d'it\'erations la boucle effectue-t-elle sur $[a,b]$ avec pas $h$~?
  \item Impl\'ementer en Python \computer\ et tester sur $f(x)=x^3-3x+1$
        avec $a=-2$, $b=2$, $h=0{,}01$.
\end{enumerate}
\end{exo}

\subsection*{\diff{3}\quad Niveau 3 \textemdash\ D\'emonstrations}

\begin{exo}[Monotonie par d\'efinition \corrige]{3}
D\'emontrer directement \`a partir de la d\'efinition ($x_1<x_2\Rightarrow f(x_1)<f(x_2)$)~:
\begin{enumerate}
  \item $f(x)=\sqrt{x}$ est croissante sur $[0,+\infty[$.
        \emph{(\'Etudier $\sqrt{x_2}-\sqrt{x_1}$ en multipliant par
        $(\sqrt{x_2}+\sqrt{x_1})/(\sqrt{x_2}+\sqrt{x_1})$.)}
  \item $g(x)=1/x$ est d\'ecroissante sur $]0,+\infty[$.
        \emph{(Calculer $g(x_2)-g(x_1)=(x_1-x_2)/(x_1 x_2)$ et \'etudier le signe.)}
\end{enumerate}
\end{exo}

\begin{exo}[Parit\'e et sym\'etries \corrige]{3}
\begin{enumerate}
  \item D\'emontrer que si $f$ est paire, sa courbe est sym\'etrique par
        rapport \`a l'axe $Oy$.
  \item D\'emontrer que la somme de deux fonctions paires est paire.
  \item Si $f$ est impaire et $f(0)$ existe, montrer que $f(0)=0$.
  \item La fonction $h(x)=f(x)+g(x)$ o\`u $f$ est paire et $g$ est impaire~:
        est-ce que $h$ est paire, impaire, ou ni l'une ni l'autre~?
\end{enumerate}
\end{exo}

\begin{exo}[$f(x)=|x^2-4|$~: \'etude compl\`ete \corrige]{3}
\begin{enumerate}
  \item Montrer que $f$ est paire.
  \item Donner l'expression de $f$ sans valeur absolue sur chaque intervalle~:
        $]-\infty,-2]$, $[-2,2]$, $[2,+\infty[$.
  \item Dresser le tableau de variations complet de $f$.
  \item Tracer la courbe.
  \item R\'esoudre $f(x)=x$ en distinguant les cas.
\end{enumerate}
\end{exo}

\begin{exo}[Signe et variation combin\'es \corrige]{3}
Pour $f(x)=\dfrac{x^2-1}{x}$ (d\'efinie sur $\R^*$)~:
\begin{enumerate}
  \item Montrer que $f$ est impaire.
  \item \'Etudier le signe de $f(x)$ sur $]-\infty,0[$ et $]0,+\infty[$.
  \item D\'emontrer que $f$ est croissante sur $]0,+\infty[$~:
        calculer $f(x_2)-f(x_1)$ pour $0<x_1<x_2$.
  \item \'Etudier les variations sur $]-\infty,0[$.
  \item Dresser le tableau de variations complet.
\end{enumerate}
\end{exo}

\begin{exo}[Encadrement et approximation \corrige]{3}
Soit $f(x)=x^2$ et $g(x)=\sqrt{x}$ sur $[0,4]$.
\begin{enumerate}
  \item Montrer que $f$ et $g$ sont inverses l'une de l'autre~:
        $f(g(x))=x$ et $g(f(x))=x$ pour $x\geq0$.
  \item En d\'eduire que les courbes de $f$ et $g$ sont sym\'etriques
        par rapport \`a la droite $y=x$.
  \item Encadrer $\sqrt{3}$ entre deux valeurs de $f$ en utilisant~:
        $f(1{,}7)=2{,}89 < 3 < f(1{,}8)=3{,}24$.
  \item Montrer que $g$ est croissante, puis que $1{,}7<\sqrt{3}<1{,}8$.
\end{enumerate}
\end{exo}

%─────────────────────────────────────────────────────────────────
\subsection*{Probl\`emes}
%─────────────────────────────────────────────────────────────────

\begin{probleme}[Optimisation~: cloture \corrigeP]{1}
Un agriculteur dispose de $100$~m de cl\^oture pour d\'elimiter un terrain
rectangulaire longeant un mur (le mur constituant un c\^ot\'e, pas besoin
de cl\^oture de ce c\^ot\'e).
\begin{enumerate}
  \item Si la largeur est $x$, exprimer la longueur en fonction de $x$.
  \item Exprimer l'aire $A(x)$ en fonction de $x$.
  \item Quel est le domaine de validit\'e de $x$~?
  \item Mettre $A(x)$ sous forme canonique et trouver la valeur de $x$
        qui maximise l'aire.
  \item Quelle est l'aire maximale~?
\end{enumerate}
\end{probleme}

\begin{probleme}[Courbe et lecture graphique \corrigeP]{1}
La courbe d'une fonction $f$ (une parabole) passe par les points
$A(-1, 3)$, $B(1,-1)$, $C(3, 3)$.
\begin{enumerate}
  \item Montrer que $f$ est de la forme $f(x)=a(x-h)^2+k$
        et trouver $h$ par sym\'etrie.
  \item Utiliser deux points pour trouver $a$ et $k$.
  \item V\'erifier avec le troisi\`eme point.
  \item Dresser le tableau de variations.
  \item R\'esoudre $f(x)=0$.
\end{enumerate}
\end{probleme}

\begin{probleme}[Optimisation~: boite \corrigeP]{2}
On d\'ecoupe des petits carr\'es de c\^ot\'e $x$ aux quatre coins d'une
feuille rectangulaire $12\times8$~cm, puis on replie pour former une bo\^ite.
\begin{enumerate}
  \item Exprimer le volume $V(x)$ en fonction de $x$.
  \item Quel est le domaine de validit\'e de $x$~?
  \item Calculer $V(1)$, $V(2)$, $V(3)$.
  \item La valeur qui maximise le volume est $x\approx1{,}57$.
        V\'erifier que $V(1{,}5)<V(1{,}57)$ et $V(2)<V(1{,}57)$.
\end{enumerate}
\end{probleme}

\begin{probleme}[Fonctions et physique~: chute libre \corrigeP]{1}
Un objet tombe depuis une hauteur de $100$~m.
Sa hauteur est $h(t)=100-5t^2$ (en m, avec $t$ en secondes).
\begin{enumerate}
  \item Quel est le domaine physique de validit\'e de $t$~?
  \item Quand l'objet touche-t-il le sol~?
  \item Dresser le tableau de variations de $h$.
  \item \`A quelle hauteur est l'objet apr\`es $3$~s~?
        Quelle vitesse repr\'esente $\dfrac{h(3)-h(0)}{3}$~?
\end{enumerate}
\end{probleme}

\begin{probleme}[Comparaison de fonctions \corrigeP]{2}
On compare la croissance de $f(x)=x^2$ et $g(x)=2x$ sur $[0,+\infty[$.
\begin{enumerate}
  \item Calculer $f(x)-g(x)$ et \'etudier son signe.
  \item Pour quelles valeurs de $x$ a-t-on $x^2\geq 2x$~?
  \item Tracer les deux courbes sur $[0,4]$.
  \item G\'en\'eraliser~: pour quelles valeurs de $a>0$ la courbe
        $y=x^2$ est-elle au-dessus de $y=ax$~?
\end{enumerate}
\end{probleme}

\begin{probleme}[Fonctions et \'economie \corrigeP]{1}
Le b\'en\'efice (en milliers d'ariary) d'une entreprise selon le nombre
$n$ de produits fabriqu\'es (en centaines) est~:
$B(n) = -n^2 + 10n - 16$.
\begin{enumerate}
  \item Quel est le domaine de validit\'e de $n$ (entier positif)~?
  \item L'entreprise fait-elle des pertes pour $n=2$~? Pour $n=8$~?
  \item Trouver les valeurs de $n$ pour lesquelles $B(n)\geq0$.
  \item Quel est le b\'en\'efice maximal et pour quel $n$~?
\end{enumerate}
\end{probleme}

\begin{probleme}[Suites et fonctions \corrigeP]{2}
On consid\`ere la suite $u_n = n^2 - 4n + 3$ pour $n\geq0$ entier.
\begin{enumerate}
  \item Calculer $u_0, u_1, u_2, u_3, u_4, u_5$.
  \item La fonction $f(x)=x^2-4x+3$ prend les m\^emes valeurs aux entiers.
        Dresser le tableau de variations de $f$ et localiser le minimum.
  \item Pour quels entiers $n$ a-t-on $u_n\leq0$~?
  \item Montrer que pour $n\geq3$, la suite est croissante.
\end{enumerate}
\end{probleme}

\begin{probleme}[Algorithme et fonctions \corrigeP]{2}
\'Ecrire un algorithme \computer\ qui~:
\begin{enumerate}
  \item Pour une liste de valeurs de $x$ en entr\'ee, calcule et affiche
        $f(x)=\sqrt{x^2+1}$.
  \item D\'etermine si $f$ semble croissante sur $[0,5]$ en
        v\'erifiant que $f(x+0{,}1)>f(x)$ pour chaque valeur.
  \item Trouve le point d'intersection de $f(x)=\sqrt{x^2+1}$
        et $g(x)=x+1$ en testant les valeurs de $x\in[0,3]$ avec un pas de $0{,}01$.
\end{enumerate}
\end{probleme}

\begin{probleme}[\'Etude de $f(x)=1/x$ \corrigeP]{2}
Pour $f(x)=1/x$, d\'efinie sur $\R^*$~:
\begin{enumerate}
  \item Montrer que $f$ est impaire.
  \item D\'emontrer que $f$ est d\'ecroissante sur $]0,+\infty[$
        en \'etudiant le signe de $f(x_2)-f(x_1)$ pour $0<x_1<x_2$.
  \item D\'emontrer que $f$ est d\'ecroissante sur $]-\infty,0[$.
  \item Pourquoi ne peut-on pas dire que $f$ est \og d\'ecroissante sur $\R^*$\fg~?
        \emph{(Comparer $f(-1)$ et $f(1)$.)}
\end{enumerate}
\end{probleme}

\begin{probleme}[Fonctions et codes postaux \corrigeP]{1}
On mod\'elise le co\^ut d'envoi d'un colis (en milliers d'ariary)
par la fonction par morceaux~:
$$C(m) = \begin{cases}2 & \text{si }0<m\leq1 \\ 3 & \text{si }1<m\leq3 \\
5 & \text{si }3<m\leq5 \\ 8 & \text{si }5<m\leq10\end{cases}$$
o\`u $m$ est la masse en kg.
\begin{enumerate}
  \item Calculer $C(0{,}5)$, $C(1)$, $C(2)$, $C(4)$, $C(7)$.
  \item Tracer la courbe de $C$ sur $[0,10]$.
  \item $C$ est-elle croissante~? Monotone~?
  \item Pour quel $m$ le co\^ut change-t-il de tranche~?
\end{enumerate}
\end{probleme}

\begin{probleme}[Maximiser une aire \corrigeP]{2}
Un rectangle est inscrit dans un demi-cercle de rayon $R=5$.
Un sommet est sur le diam\`etre, les deux autres sur le demi-cercle.
Si le sommet sur le diam\`etre est en $(-x,0)$ et $(x,0)$
($0<x<5$), la hauteur du rectangle est $\sqrt{25-x^2}$.
\begin{enumerate}
  \item Exprimer l'aire $A(x)$ du rectangle en fonction de $x$.
  \item Calculer $A(1)$, $A(2)$, $A(3)$, $A(4)$.
  \item L'aire maximale semble-t-elle atteinte pour $x=\sqrt{2}\cdot5/2\approx3{,}54$~?
        V\'erifier num\'eriquement.
  \item Quel est le p\'erim\`etre du rectangle correspondant~?
\end{enumerate}
\end{probleme}

\begin{probleme}[Double transformation \corrigeP]{2}
On part de $f(x)=x^2$.
\begin{enumerate}
  \item \'Ecrire la fonction obtenue par~:
        translation de $3$ vers la droite, puis sym\'etrie par rapport \`a $Ox$,
        puis translation de $2$ vers le haut.
  \item Trouver le sommet de cette parabole.
  \item La m\^eme parabole en changeant l'ordre~:
        sym\'etrie par $Ox$, puis translation de $3$ \`a droite, puis $2$ vers le haut.
        Obtient-on la m\^eme fonction~?
\end{enumerate}
\end{probleme}

%─────────────────────────────────────────────────────────────────
\subsection*{D\'efis}
%─────────────────────────────────────────────────────────────────

\begin{challenge}[Vrai ou faux~: monotonie et image]
Pour chaque affirmation, dire si elle est vraie ou fausse et justifier~:
\begin{enumerate}
  \item \og Si $f$ est croissante et $f(3)>0$, alors $f(x)>0$ pour tout $x>3$.\fg
  \item \og Si $f(a)=f(b)$ avec $a\ne b$, alors $f$ n'est pas monotone.\fg
  \item \og Si $f$ est croissante et $g$ est croissante, alors $f+g$ est croissante.\fg
  \item \og Si $f$ est croissante, alors $f\circ f$ est croissante.\fg
  \item \og Une fonction paire ne peut pas \^etre strictement croissante sur $\R$.\fg
\end{enumerate}
\end{challenge}

\begin{challenge}[Minimum de $f(x)=x+1/x$]
Pour $f(x)=x+\dfrac{1}{x}$ sur $]0,+\infty[$~:
\begin{enumerate}
  \item Calculer $f(1/2)$, $f(1)$, $f(2)$, $f(3)$.
  \item Montrer que $f(x)\geq2$ pour tout $x>0$.
        \emph{(Utiliser $f(x)-2=x+\frac{1}{x}-2=\frac{x^2-2x+1}{x}=\frac{(x-1)^2}{x}$.)}
  \item Quand l'\'egalit\'e $f(x)=2$ est-elle atteinte~?
  \item De m\^eme, montrer que pour $x<0$, $f(x)\leq-2$.
\end{enumerate}
\end{challenge}

\begin{challenge}[Graphe et transformations enchaîn\'ees]
Soit $f(x)=|x|$. On applique les transformations successives~:
\begin{enumerate}
  \item \'Ecrire la courbe de $g(x)=|x-2|+1$ et identifier
        la translation appliqu\'ee \`a $f$.
  \item \'Ecrire $h(x)=-|x+1|+3$ et tracer sa courbe.
  \item Trouver toutes les fonctions de la forme $a|x-h|+k$
        ($a\ne0$) qui passent par les points $(0,2)$ et $(4,2)$.
  \item Les courbes de $f(x)=|x|$ et $g(x)=|x-2|$ se coupent-elles~?
        Trouver le(s) point(s) d'intersection.
\end{enumerate}
\end{challenge}

\begin{challenge}[Fonction constante par morceaux]
La \textbf{partie enti\`ere} $\lfloor x\rfloor$ est le plus grand entier $\leq x$.
\begin{enumerate}
  \item Calculer $\lfloor 2{,}7\rfloor$, $\lfloor -1{,}3\rfloor$,
        $\lfloor 3\rfloor$, $\lfloor -2\rfloor$.
  \item Tracer le graphe de $f(x)=\lfloor x\rfloor$ sur $[-2,3]$.
  \item Pour $g(x)=x-\lfloor x\rfloor$ (partie fractionnaire)~:
        calculer $g(1{,}7)$, $g(-0{,}3)$, $g(5)$.
        Tracer $g$ sur $[-1,3]$.
  \item Sur quel intervalle $g(x)$ est-elle croissante~?
        Est-elle continue~?
\end{enumerate}
\end{challenge}

\begin{challenge}[Fonctions et suite de nombres]
Soit $f(x)=x^2-2$ et la suite d\'efinie par $u_0=2$ et $u_{n+1}=f(u_n)$.
\begin{enumerate}
  \item Calculer $u_1$, $u_2$, $u_3$, $u_4$.
  \item Cette suite est-elle croissante, d\'ecroissante, non monotone~?
  \item La suite est-elle minor\'ee~? Major\'ee~?
  \item Comparer avec la suite $v_n=f^n(1)$ (en partant de $1$).
        Que semble-t-il se passer~?
\end{enumerate}
\end{challenge}

\begin{challenge}[Sch\'ema de point fixe]
Un \textbf{point fixe} de $f$ est un r\'eel $c$ tel que $f(c)=c$,
c'est-\`a-dire un point d'intersection de $y=f(x)$ et $y=x$.
\begin{enumerate}
  \item Trouver les points fixes de $f(x)=x^2-x+1$.
        \emph{(R\'esoudre $x^2-x+1=x$.)}
  \item Trouver les points fixes de $g(x)=\sqrt{x}$ ($x\geq0$).
  \item Trouver les points fixes de $h(x)=\dfrac{x+3}{2}$.
  \item La suite $u_{n+1}=h(u_n)$ avec $u_0=0$~: calculer $u_1,\ldots,u_6$.
        Que semble-t-elle faire~?
\end{enumerate}
\end{challenge}

%─────────────────────────────────────────────────────────────────
\begin{bilan}
\begin{itemize}
  \item \textbf{Domaine}~: exclure les z\'eros du d\'enominateur
        et les valeurs rendant un radical n\'egatif.
  \item \textbf{Image~:} $f(x_0)$. \textbf{Ant\'ec\'edent~:} r\'esoudre $f(x)=y_0$.
  \item \textbf{Tableau de variations~:} fl\`eches $\nearrow$ (croissant)
        et $\searrow$ (d\'ecroissant)~; valeurs aux bornes et extremums.
  \item \textbf{Fonctions de r\'ef\'erence~:}
        $x^2$ (paire, min en 0)~; $x^3$ (impaire, croissante)~;
        $\sqrt{x}$ (croissante sur $[0,+\infty[$)~;
        $|x|$ (paire, min en 0)~; $1/x$ (impaire, d\'ecr. sur $\R^*$).
  \item \textbf{Transformations~:} $f(x)+k$ (translation vert.)~;
        $f(x+k)$ (translation horiz.)~; $-f(x)$ (sym.\ $Ox$)~;
        $f(-x)$ (sym.\ $Oy$).
\end{itemize}

%─────────────────────────────────────────────────────────────────
\subsection*{Logique \& Algorithmes}
%─────────────────────────────────────────────────────────────────

\begin{exo}[\logiq\ Logique et domaine de d\'efinition \corrige]{1}
\'Enoncer sous forme logique (avec $\exists$ ou $\forall$)~:
\begin{enumerate}
  \item Le domaine de d\'efinition de $f(x)=\sqrt{x-2}$.
  \item La condition pour que $g(x)=1/(x^2-4)$ soit d\'efinie.
  \item La proposition~: \og toute valeur de $f$ est positive\fg.
  \item La n\'egation de~: \og $f$ admet un ant\'ec\'edent pour toute valeur\fg.
\end{enumerate}
\end{exo}

\begin{exo}[\logiq\ Implication et croissance \corrige]{2}
\begin{enumerate}
  \item \og $f$ croissante sur $[a,b]$ et $x_1<x_2$ dans $[a,b]$
        $\Rightarrow f(x_1)<f(x_2)$\fg~: c'est la d\'efinition.
        \'Enoncer la contrapos\'ee.
  \item Si $f(3)=5$ et $f(5)=3$, que peut-on conclure sur la monotonie de $f$~?
        Quel raisonnement utilise-t-on~?
  \item D\'emontrer que $f(x)=x^2$ n'est pas monotone sur $\R$.
        Donner un contre-exemple.
\end{enumerate}
\end{exo}

\begin{exo}[\logiq\ Conditions sur les extremums \corrige]{1}
Pour une fonction $f$ de tableau de variations donn\'e~:
\begin{enumerate}
  \item \'Enoncer~: \og $f$ admet un minimum absolu\fg\ avec des quantificateurs.
  \item La proposition~: \og si $f(a)=0$, alors $a$ est un minimum\fg.
        Vraie ou fausse~? Contre-exemple si fausse.
  \item \og $f$ est born\'ee si et seulement si elle admet un minimum
        et un maximum\fg~: est-ce exact~?
\end{enumerate}
\end{exo}

\begin{exo}[\logiq\ N\'egation de propositions fonctionnelles \corrige]{2}
\'Enoncer la n\'egation et d\'eterminer laquelle est vraie~:
\begin{enumerate}
  \item $\forall x\in\R,\ x^2\geq0$.
  \item $\exists x\in\R^*,\ 1/x=x$.
  \item $\forall x>0,\ \sqrt{x}<x$.
  \item $\forall x\in[-1,1],\ |x|\leq1$.
\end{enumerate}
\end{exo}

\begin{exo}[\logiq\ Parit\'e~: d\'emonstrations logiques \corrige]{2}
\begin{enumerate}
  \item La somme de deux fonctions paires est-elle toujours paire~?
        D\'emontrer ou donner un contre-exemple.
  \item Le produit d'une fonction paire et d'une fonction impaire~:
        quelle est sa parit\'e~? D\'emontrer.
  \item Une fonction peut-elle \^etre \`a la fois paire et impaire~?
        Si oui, laquelle~?
  \item Formuler les r\'esultats pr\'ec\'edents comme des implications.
\end{enumerate}
\end{exo}

\begin{exo}[\logiq\ Logique et transformations de courbes \corrige]{2}
\begin{enumerate}
  \item La proposition~: \og si $f$ est croissante, alors $-f$ est d\'ecroissante\fg~:
        d\'emontrer.
  \item \og Si $f$ est paire, alors $f(x)+f(-x)=2f(x)$ pour tout $x$\fg~:
        vraie ou fausse~?
  \item D\'emontrer~: si $f$ est born\'ee par $M$, alors $|f(x)|\leq M$.
        Formuler comme implication.
\end{enumerate}
\end{exo}

\begin{exo}[\algo\ Algorithme~: tableau de valeurs d'une fonction \computer \corrige]{1}
\'Ecrire un programme Python qui calcule et affiche le tableau de valeurs
de $f(x)=x^2-3x+1$ pour $x\in\{-2,-1,0,1,2,3,4\}$.
Puis d\'eterminer graphiquement les z\'eros de $f$.
V\'erifier avec la formule du discriminant.
\end{exo}

\begin{exo}[\algo\ Algorithme~: trouver le minimum par balayage \computer \corrige]{2}
\'Ecrire un algorithme qui approche le minimum de $f(x)=x^2-4x+1$
sur $[-1,5]$ en balayant avec un pas de $0{,}01$.
Afficher la valeur approch\'ee et la comparer avec le minimum exact.
\end{exo}

\begin{exo}[\algo\ Algorithme~: d\'etecter les z\'eros par changement de signe \computer \corrige]{2}
(Voir exercice \logicoalgo\ dans ce chapitre.)
Tester sur~:
\begin{enumerate}
  \item $f(x)=x^3-3x+1$ sur $[-2,2]$ avec $h=0{,}001$.
  \item $g(x)=|x|-1$ sur $[-3,3]$.
  \item $h(x)=\sqrt{x}-x$ sur $[0,2]$.
\end{enumerate}
Comparer avec les racines exactes.
\end{exo}

\begin{exo}[\algo\ Algorithme~: v\'erifier la monotonie \computer \corrige]{2}
\'Ecrire un programme Python qui~:
\begin{enumerate}
  \item V\'erifie si $f$ est croissante sur $[a,b]$ en testant
        $f(x_i)\leq f(x_{i+1})$ pour un balayage avec pas $h$.
  \item Teste $f(x)=x^2$ sur $[0,3]$ et $[-3,0]$.
  \item Teste $g(x)=x^3-x$ sur $[-2,2]$~: trouver les sous-intervalles
        de monotonie.
\end{enumerate}
\end{exo}

\begin{exo}[\algo\ Algorithme~: point fixe par it\'eration \computer \corrige]{3}
\'Ecrire un programme qui calcule la suite $u_{n+1}=h(u_n)$
avec $h(x)=(x+3)/2$, $u_0=0$, jusqu'\`a ce que $|u_{n+1}-u_n|<10^{-6}$.
\begin{enumerate}
  \item Afficher chaque terme et le nombre d'it\'erations.
  \item V\'erifier que la limite est le point fixe de $h$.
  \item Tester avec d'autres valeurs de $u_0$.
\end{enumerate}
\end{exo}

\begin{exo}[\algo\ Algorithme~: maximum d'une liste de valeurs \computer \corrige]{1}
\'Ecrire un programme qui~:
\begin{enumerate}
  \item Re\c coit une liste de valeurs de $f$ en $n$ points.
  \item Retourne le maximum, le minimum et leur position.
  \item Tester sur $f(x)=\sin(x)$ \'evalu\'e en $100$ points de $[0,2\pi]$.
\end{enumerate}
\end{exo}

%─────────────────────────────────────────────────────────────────
\subsection*{Probl\`emes Logique \& Algorithmes}
%─────────────────────────────────────────────────────────────────

\begin{probleme}[\logiq\ Logique et optimisation \corrigeP]{2}
\begin{enumerate}
  \item Montrer~: \og $f$ admet un maximum en $x_0$\fg\ implique
        \og $f(x_0)\geq f(x)$ pour tout $x$ dans son domaine\fg.
        Formuler l'\'equivalence exacte.
  \item Pour $f(x)=-x^2+4x-3$~: d\'emontrer logiquement que le maximum
        est atteint en $x=2$ en montrant que $f(x)\leq f(2)$ pour tout $x$.
  \item \'Enoncer la condition exacte pour qu'un trin\^ome $ax^2+bx+c$
        admette un minimum~: condition sur $a$ (CN, CS, CNS~?).
\end{enumerate}
\end{probleme}

\begin{probleme}[\logiq\ Propri\'et\'es logiques de la composition \corrigeP]{2}
\begin{enumerate}
  \item Montrer que la compos\'ee de deux fonctions croissantes est croissante.
        Formuler l'implication et la d\'emontrer.
  \item La compos\'ee de deux fonctions impaires est-elle paire ou impaire~?
        D\'emontrer.
  \item \'Enoncer la contrapos\'ee de~:
        \og Si $f$ est strictement croissante et $g$ est strictement croissante,
        alors $f\circ g$ est strictement croissante.\fg
        La contrapos\'ee est-elle utile~?
\end{enumerate}
\end{probleme}

\begin{probleme}[\algo\ Algorithme~: \'etude compl\`ete d'une fonction \computer \corrigeP]{2}
\'Ecrire un programme Python complet qui, pour une fonction $f$ donn\'ee~:
\begin{enumerate}
  \item Calcule le tableau de valeurs sur $[a,b]$.
  \item D\'etermine (approximativement) les z\'eros, le maximum et le minimum.
  \item Trace la courbe avec \texttt{matplotlib}.
  \item Affiche un r\'esum\'e des propri\'et\'es.
\end{enumerate}
Tester sur $f(x)=x^3-3x+1$.
\end{probleme}

\begin{probleme}[\algo\ Algorithme~: optimisation de l'aire d'un rectangle \computer \corrigeP]{2}
(Voir probl\`eme de cl\^oture.)
Un rectangle de p\'erim\`etre $P=20$ a ses c\^ot\'es $x$ et $10-x$.
\'Ecrire un programme Python qui~:
\begin{enumerate}
  \item Calcule $A(x)=(10-x)\times x$ pour $x\in[0,10]$ avec pas $0{,}001$.
  \item Trouve le $x$ qui maximise $A$.
  \item Compare avec la solution exacte.
  \item Trace $A(x)$.
\end{enumerate}
\end{probleme}

%─────────────────────────────────────────────────────────────────
\subsection*{D\'efis Logique \& Algorithmes}
%─────────────────────────────────────────────────────────────────

\begin{challenge}[\logiq\ Th\'eor\`eme des valeurs interm\'ediaires]
On admet que si $f$ est continue sur $[a,b]$ et $f(a)<0<f(b)$,
alors $\exists c\in]a,b[$ tel que $f(c)=0$.
\begin{enumerate}
  \item Formuler ce th\'eor\`eme avec des quantificateurs.
  \item \'Enoncer la contrapos\'ee.
  \item Utiliser le TVI pour montrer que $x^3-2=0$ a une solution dans $]1,2[$.
  \item \'Enoncer la r\'eciproque du TVI~: est-elle vraie~?
        Donner un exemple ou contre-exemple.
\end{enumerate}
\end{challenge}

\begin{challenge}[\algo\ Algorithme de bissection g\'en\'eralis\'e \computer]
\'Ecrire un programme Python g\'en\'eral de bissection (dichotomie)
qui trouve une racine de toute fonction continue $f$ sur $[a,b]$
avec $f(a)f(b)<0$, \`a une pr\'ecision $\varepsilon$ donn\'ee.
\begin{enumerate}
  \item Impl\'ementer l'algorithme.
  \item Afficher le nombre d'it\'erations n\'ecessaires.
  \item Tester sur $f(x)=x^3-2$ (trouver $\sqrt[3]{2}$).
  \item Tester sur $f(x)=\cos(x)-x$ sur $[0,\pi/2]$.
\end{enumerate}
\end{challenge}

\end{bilan}
